\subsection{Verarbeitung und Nutzung der Wetterdaten}Allgemein ist die Integration externer Datenquellen mit Herausforderungen verbunden, insbesondere hinsichtlich Verfügbarkeit und Zuverlässigkeit der bereitgestellten Dienste~\cite{7123563}. 
Bei Ausfall der Verbindung oder des Dienstes können keine aktuellen Daten abgerufen werden. 
Daher ist eine geeignete Fehlerbehandlung erforderlich, um weiterhin einen stabilen Systembetrieb sicherzustellen.

Für die Darstellung der Wetterdaten wird eine vereinfachte Klassifikation der Messwerte verwendet. 
Kontinuierliche Größen wie die Temperatur werden in diskrete Bereiche unterteilt. 
Dadurch wird eine schwellwertabhängige Darstellung im Fall der Wortuhr ermöglicht. 
Beispielsweise können Temperaturbereiche in Kategorien wie kühl, frisch oder warm eingeteilt werden.

Neben der Temperatur können auch Wetterzustände wie ein klarer Himmel, Bewölkung oder Niederschlag berücksichtigt werden. 
Diese Zustände werden aus den bereitgestellten Daten abgeleitet und ebenfalls über Schwellenwerte klassifiziert. 
Dadurch wird eine schnelle Interpretation und Darstellung der aktuellen Wetterlage ermöglicht. 

Die Verarbeitung der Daten erfolgt zyklisch, wobei in definierten Zeitintervallen neue Wetterinformationen vom externen Dienst abgerufen werden. 
Zwischen den Aktualisierungen werden die zuletzt gültigen Werte weiterhin verwendet, um eine kontinuierliche Anzeige sicherzustellen.
Durch die Nutzung standardisierter Schnittstellen können Datenquellen einfach integriert und in bestehende Systeme eingebunden werden.
